% Tapis de Quillen, with Quillen's offprint and preprint (cote 162-5, typed part): archive cote 162-5, site page 41 (PDF page 42).
% typescript. Draft compared: none.
% AI-transcribed by Claude Opus 5.5 (claude-opus-5-5, adaptive thinking, effort high), Message Batches API, batch pilotnext (msgbatch_01A9sdu8SYq23bEtqX9tUncQ), 2026-11-13; not yet checked by eye.
\scanpage{162-5/41}
qu'il conviendra de poser les définitions adéquates). - Peut-être Quillen
\margin{oui} \margin{$\|$}a-t-il une définition semi-simpliciale plus directe, sans passer par des
n-catégories, des $K_i$ et plus précisément du H-espace qui devrait correspon-%
dre à C. \xout{Dans le cas} Ce n'est certainement pas \inserted[typed]{en général} l'ensemble semi-simplicial
SC \xout{engendré à} associé à C, puisque dans le cas d'une catégorie additive,
C est homotopiquement trivial (en effet le foncteur identique est ``homotope
au foncteur nul, puisqu'il suffit d'avoir un homomorphisme d'un foncteur
dans l'autre; d'où une équivalence d'homotopie entre C et la catégorie
nulle). Pourtant, il semble bien que dans le cas où C est une Gr-catégorie
\margin{\unsure{$\xi$|$\zeta$}} \margin{$|$}i.e. admet des inverses, alors ce soit bien la construction ``naïve'' SC qui
est la bonne, puisqu'on trouve bien un $\hand{\pi}_0$ et un $\hand{\pi}_1$ qui sont aussi les
$\hand{\pi}_0$ et $\hand{\pi}_1$ au sens  catégorie sans Gr-structure.

\margin{À expliquer : pour des Gr-n-catégories, les $\pi_i$ au sens
\unsure{$\xi$|$\zeta$} $|$ semi-simplicial sont aussi les $\pi_i$ au sens Gr-n-catégories,
savoir les objets \unsure{:|;} isom pris dans divers $\underline{\mathrm{Hom}}(1,1)$ ...... \deleted{Il y a
aussi une interprétation de l'opération \unsure{$\Omega$ : espace des lacets}
en termes d'opération \unsure{de|du} \unsure{shift} géométrique évidente sur les
n-catégories .....}}

\mspage{41}
